%% ============================================================
%%  Aula 6 — Gradientes de política II: variância, confiança e PPO
%%  Adaptação em pt-BR do CS234 (Stanford, Emma Brunskill), Lecture 6
%%  Seção de gradientes avançados baseada nos slides de Joshua Achiam
%%  Compilar: latexmk -xelatex aula06.tex   (duas passadas)
%%  Figuras de simulação: fig-variancia.tex e fig-variancia-treino.tex,
%%  geradas por sim/simulacao.py e sim/figuras.py (ver anexo).
%% ============================================================
\documentclass[aspectratio=169, 11pt]{beamer}
%% .sty do projeto na raiz do repositório (../)
\makeatletter
\def\input@path{{./}{../}}
\makeatother

\usetheme{AprendizadoRL}      % ANTES do babel — fontspec precisa vir primeiro
\usepackage{babel}
\babelprovide[import, main]{brazilian}
\usepackage{booktabs}
\usepackage{colortbl}
\usepackage{pgfplots}
\pgfplotsset{compat=1.18}

\title[Aula 6 — Gradientes de política II]{Aula 6: Gradientes de política II}
\subtitle[Aprendizagem por Reforço]{Linha de base \textbullet\ Ator--crítico \textbullet\ Limite de desempenho relativo \textbullet\ PPO}
\author[Prof. Marcelo K. Albertini]{Prof. Marcelo Keese Albertini \\
  \footnotesize Universidade Federal de Uberlândia \\
  \footnotesize adaptado de Emma Brunskill, CS234 (Stanford); seção avançada de Joshua Achiam}
\institute{Universidade Federal de Uberlândia}
\date{2026}

% ---- Macros locais da aula
\newcommand{\ghat}{\hat{g}}
\newcommand{\Ahat}{\hat{A}}
\newcommand{\KL}{D_{\mathrm{KL}}}
\newcommand{\KLbar}{\overline{D}_{\mathrm{KL}}}
\newcommand{\Var}{\operatorname{Var}}
\DeclareMathOperator{\clipop}{clip}
\newcommand{\azul}[1]{{\color{rlAzulClaro}#1}}

\begin{document}

\begin{frame}[plain, noframenumbering]\titlepage\end{frame}

%% ============================================================
\begin{frame}{Uma atualização, dois gargalos}
  \framesubtitle{Maximizamos o retorno esperado com fator de desconto $\gamma = 1$}

  \[
    J(\theta) \defeq \E_{\tau\sim\pol_\theta}\Big[\sum_{t=0}^{\infty}\gamma^t r_t\Big]
    \qquad
    G_t \defeq \sum_{t'=t}^{\infty}\gamma^{t'-t}\, r_{t'}
    \qquad
    \theta_{k+1} = \theta_k + \alpha\,\ghat_k
  \]

  \begin{columns}[T]
    \begin{column}{0.48\textwidth}
      \begin{teobox}{I · Estatística: quão precisa é a direção $\ghat_k$?}
        \small O REINFORCE é não-viesado, mas sua variância é tão alta que
        $\ghat_k$ costuma apontar para o lado errado.\\[0.3em]
        \textbf{Rota:} linha de base $\to$ crítico de valor $\to$ vantagem.
      \end{teobox}
    \end{column}
    \begin{column}{0.48\textwidth}
      \begin{exbox}[title={II · Geometria: qual o maior passo seguro?}]{}
        \small Dados \emph{on-policy} expiram após um passo, e medir
        $\|\Delta\theta\|$ não controla a mudança real na política.\\[0.3em]
        \textbf{Rota:} métrica KL/Fisher $\to$ objetivo substituto $\to$ PPO.
      \end{exbox}
    \end{column}
  \end{columns}
\end{frame}

%% ============================================================
\begin{frame}{Algoritmo REINFORCE}
  \framesubtitle{Recordando a Aula 5, com a causalidade temporal: $G_t$ no lugar de $R(\tau)$, para $m$ trajetórias amostradas}

  \[
    \ghat = \frac{1}{m}\sum_{i=1}^{m}\sum_{t=0}^{T-1}\gamma^t\,
      \nabla_\theta\log\pol_\theta\big(a_t^{(i)}\mid s_t^{(i)}\big)\,G_t^{(i)}
  \]

  \begin{alertabox}[title={Não-viesado não significa utilizável}]
    Com um orçamento finito de $m$ amostras, quem governa a qualidade do passo
    é o erro quadrático médio:
    \[
      \mathrm{EQM}(\ghat) = \big\|\,\text{viés}\,\big\|^2 + \frac{\Var}{m}
    \]
    Variância enorme produz passos quase ortogonais --- ou até opostos --- a $\nabla J(\theta)$.
  \end{alertabox}
\end{frame}

%% ============================================================
\begin{frame}{O que significa a variância de um estimador}
  \begin{center}
    % Gerado por figs.py (mesma semente do deck). Painéis de variância alta e baixa.
\begin{tikzpicture}
\clip (0,0) rectangle (6.240,3.744);
\fill[rlFundo] (0,0) rectangle (6.240,3.744);
\fill[rlLaranja!14] (8.303,-12.783) -- (-5.651,15.123) -- (-19.604,8.147) -- (-5.651,-19.760) -- cycle;
\draw[rlLaranja, dashed] (8.303,-12.783) -- (-5.651,15.123);
\draw[rlLaranja, opacity=0.25] (1.326,1.170) -- (2.654,3.382);
\draw[rlLaranja, opacity=0.25] (1.326,1.170) -- (3.571,1.711);
\draw[rlLaranja, opacity=0.25] (1.326,1.170) -- (1.460,1.686);
\draw[rlLaranja, opacity=0.25] (1.326,1.170) -- (3.604,-4.022);
\draw[rlLaranja, opacity=0.25] (1.326,1.170) -- (2.840,1.856);
\draw[rlLaranja, opacity=0.25] (1.326,1.170) -- (-3.147,-1.430);
\draw[rlLaranja, opacity=0.25] (1.326,1.170) -- (1.024,2.744);
\draw[rlLaranja, opacity=0.25] (1.326,1.170) -- (1.921,4.506);
\draw[rlLaranja, opacity=0.25] (1.326,1.170) -- (4.287,2.316);
\draw[rlLaranja, opacity=0.25] (1.326,1.170) -- (5.064,-0.787);
\draw[rlLaranja, opacity=0.25] (1.326,1.170) -- (1.425,-2.231);
\draw[rlLaranja, opacity=0.25] (1.326,1.170) -- (-1.033,1.558);
\draw[rlLaranja, opacity=0.25] (1.326,1.170) -- (6.723,0.729);
\draw[rlLaranja, opacity=0.25] (1.326,1.170) -- (0.085,0.559);
\draw[rlLaranja, opacity=0.25] (1.326,1.170) -- (2.329,8.826);
\draw[rlLaranja, opacity=0.25] (1.326,1.170) -- (6.954,3.767);
\draw[rlLaranja, opacity=0.25] (1.326,1.170) -- (1.104,4.220);
\draw[rlLaranja, opacity=0.25] (1.326,1.170) -- (4.620,2.736);
\draw[rlLaranja, opacity=0.25] (1.326,1.170) -- (4.983,4.639);
\draw[rlLaranja, opacity=0.25] (1.326,1.170) -- (3.343,1.642);
\draw[rlLaranja, opacity=0.25] (1.326,1.170) -- (3.957,5.953);
\draw[rlLaranja, opacity=0.25] (1.326,1.170) -- (3.697,2.568);
\draw[rlLaranja, opacity=0.25] (1.326,1.170) -- (4.844,0.685);
\draw[rlLaranja, opacity=0.25] (1.326,1.170) -- (4.705,-0.421);
\draw[rlLaranja, opacity=0.25] (1.326,1.170) -- (3.198,3.308);
\draw[rlLaranja, opacity=0.25] (1.326,1.170) -- (2.671,2.295);
\draw[rlLaranja, opacity=0.25] (1.326,1.170) -- (2.155,3.503);
\draw[rlLaranja, opacity=0.25] (1.326,1.170) -- (7.091,3.045);
\draw[rlLaranja, opacity=0.25] (1.326,1.170) -- (7.000,2.698);
\draw[rlLaranja, opacity=0.25] (1.326,1.170) -- (3.365,1.350);
\fill[rlLaranja, opacity=0.75] (2.654,3.382) circle (1.3pt);
\fill[rlLaranja, opacity=0.75] (3.571,1.711) circle (1.3pt);
\fill[rlLaranja, opacity=0.75] (1.460,1.686) circle (1.3pt);
\fill[rlLaranja, opacity=0.75] (3.604,-4.022) circle (1.3pt);
\fill[rlLaranja, opacity=0.75] (2.840,1.856) circle (1.3pt);
\fill[rlLaranja, opacity=0.75] (-3.147,-1.430) circle (1.3pt);
\fill[rlLaranja, opacity=0.75] (1.024,2.744) circle (1.3pt);
\fill[rlLaranja, opacity=0.75] (1.921,4.506) circle (1.3pt);
\fill[rlLaranja, opacity=0.75] (4.287,2.316) circle (1.3pt);
\fill[rlLaranja, opacity=0.75] (5.064,-0.787) circle (1.3pt);
\fill[rlLaranja, opacity=0.75] (1.425,-2.231) circle (1.3pt);
\fill[rlLaranja, opacity=0.75] (-1.033,1.558) circle (1.3pt);
\fill[rlLaranja, opacity=0.75] (6.723,0.729) circle (1.3pt);
\fill[rlLaranja, opacity=0.75] (0.085,0.559) circle (1.3pt);
\fill[rlLaranja, opacity=0.75] (2.329,8.826) circle (1.3pt);
\fill[rlLaranja, opacity=0.75] (6.954,3.767) circle (1.3pt);
\fill[rlLaranja, opacity=0.75] (1.104,4.220) circle (1.3pt);
\fill[rlLaranja, opacity=0.75] (4.620,2.736) circle (1.3pt);
\fill[rlLaranja, opacity=0.75] (4.983,4.639) circle (1.3pt);
\fill[rlLaranja, opacity=0.75] (3.343,1.642) circle (1.3pt);
\fill[rlLaranja, opacity=0.75] (3.957,5.953) circle (1.3pt);
\fill[rlLaranja, opacity=0.75] (3.697,2.568) circle (1.3pt);
\fill[rlLaranja, opacity=0.75] (4.844,0.685) circle (1.3pt);
\fill[rlLaranja, opacity=0.75] (4.705,-0.421) circle (1.3pt);
\fill[rlLaranja, opacity=0.75] (3.198,3.308) circle (1.3pt);
\fill[rlLaranja, opacity=0.75] (2.671,2.295) circle (1.3pt);
\fill[rlLaranja, opacity=0.75] (2.155,3.503) circle (1.3pt);
\fill[rlLaranja, opacity=0.75] (7.091,3.045) circle (1.3pt);
\fill[rlLaranja, opacity=0.75] (7.000,2.698) circle (1.3pt);
\fill[rlLaranja, opacity=0.75] (3.365,1.350) circle (1.3pt);
\fill[rlLaranja, opacity=0.75] (3.518,4.737) circle (1.3pt);
\fill[rlLaranja, opacity=0.75] (0.216,3.295) circle (1.3pt);
\fill[rlLaranja, opacity=0.75] (0.487,1.537) circle (1.3pt);
\fill[rlLaranja, opacity=0.75] (1.695,1.847) circle (1.3pt);
\fill[rlLaranja, opacity=0.75] (2.255,3.355) circle (1.3pt);
\fill[rlLaranja, opacity=0.75] (5.530,1.991) circle (1.3pt);
\fill[rlLaranja, opacity=0.75] (-0.473,3.939) circle (1.3pt);
\fill[rlLaranja, opacity=0.75] (4.583,-0.497) circle (1.3pt);
\fill[rlLaranja, opacity=0.75] (2.278,2.532) circle (1.3pt);
\fill[rlLaranja, opacity=0.75] (0.499,2.391) circle (1.3pt);
\fill[rlLaranja, opacity=0.75] (2.538,-0.374) circle (1.3pt);
\fill[rlLaranja, opacity=0.75] (2.239,3.893) circle (1.3pt);
\fill[rlLaranja, opacity=0.75] (6.292,6.007) circle (1.3pt);
\fill[rlLaranja, opacity=0.75] (3.037,-1.648) circle (1.3pt);
\fill[rlLaranja, opacity=0.75] (0.246,5.263) circle (1.3pt);
\fill[rlLaranja, opacity=0.75] (7.695,0.889) circle (1.3pt);
\fill[rlLaranja, opacity=0.75] (0.718,1.119) circle (1.3pt);
\fill[rlLaranja, opacity=0.75] (1.281,0.943) circle (1.3pt);
\fill[rlLaranja, opacity=0.75] (1.680,2.817) circle (1.3pt);
\fill[rlLaranja, opacity=0.75] (3.318,2.650) circle (1.3pt);
\fill[rlLaranja, opacity=0.75] (4.272,-0.010) circle (1.3pt);
\fill[rlLaranja, opacity=0.75] (6.005,3.835) circle (1.3pt);
\fill[rlLaranja, opacity=0.75] (2.455,1.826) circle (1.3pt);
\fill[rlLaranja, opacity=0.75] (3.513,0.564) circle (1.3pt);
\fill[rlLaranja, opacity=0.75] (2.174,-1.510) circle (1.3pt);
\fill[rlLaranja, opacity=0.75] (5.058,1.846) circle (1.3pt);
\fill[rlLaranja, opacity=0.75] (3.590,2.320) circle (1.3pt);
\fill[rlLaranja, opacity=0.75] (0.736,-1.525) circle (1.3pt);
\fill[rlLaranja, opacity=0.75] (1.749,1.329) circle (1.3pt);
\fill[rlLaranja, opacity=0.75] (2.269,0.994) circle (1.3pt);
\fill[rlLaranja, opacity=0.75] (5.572,-0.367) circle (1.3pt);
\fill[rlLaranja, opacity=0.75] (0.491,2.999) circle (1.3pt);
\fill[rlLaranja, opacity=0.75] (-2.588,1.313) circle (1.3pt);
\fill[rlLaranja, opacity=0.75] (4.857,2.866) circle (1.3pt);
\fill[rlLaranja, opacity=0.75] (4.039,-0.328) circle (1.3pt);
\fill[rlLaranja, opacity=0.75] (7.505,3.431) circle (1.3pt);
\fill[rlLaranja, opacity=0.75] (1.523,4.668) circle (1.3pt);
\fill[rlLaranja, opacity=0.75] (-0.784,-1.960) circle (1.3pt);
\fill[rlLaranja, opacity=0.75] (0.463,3.208) circle (1.3pt);
\fill[rlLaranja, opacity=0.75] (-1.485,-0.233) circle (1.3pt);
\draw[-{Stealth[length=6pt]}, line width=1.6pt, rlAzul] (1.326,1.170) -- (3.354,2.184);
\fill[rlAzul] (1.326,1.170) circle (1.8pt);
\node[anchor=north west, font=\scriptsize\bfseries, text=rlLaranja] at (0.08,3.684) {Variância alta};
\node[font=\scriptsize, text=rlAzul] at (2.574,1.365) {$\nabla J(\theta)$};
\node[anchor=south west, font=\tiny, text=rlLaranja] at (0.05,0.03) {região onde o passo piora o objetivo};
\draw[rlCinza!40] (0,0) rectangle (6.240,3.744);
\end{tikzpicture}
\hfill
\begin{tikzpicture}
\clip (0,0) rectangle (6.240,3.744);
\fill[rlFundo] (0,0) rectangle (6.240,3.744);
\fill[rlLaranja!14] (8.303,-12.783) -- (-5.651,15.123) -- (-19.604,8.147) -- (-5.651,-19.760) -- cycle;
\draw[rlLaranja, dashed] (8.303,-12.783) -- (-5.651,15.123);
\draw[rlAzulClaro, opacity=0.25] (1.326,1.170) -- (3.285,2.400);
\draw[rlAzulClaro, opacity=0.25] (1.326,1.170) -- (3.179,2.253);
\draw[rlAzulClaro, opacity=0.25] (1.326,1.170) -- (3.315,1.861);
\draw[rlAzulClaro, opacity=0.25] (1.326,1.170) -- (3.316,2.651);
\draw[rlAzulClaro, opacity=0.25] (1.326,1.170) -- (2.923,2.269);
\draw[rlAzulClaro, opacity=0.25] (1.326,1.170) -- (3.362,2.408);
\draw[rlAzulClaro, opacity=0.25] (1.326,1.170) -- (3.515,1.746);
\draw[rlAzulClaro, opacity=0.25] (1.326,1.170) -- (3.108,1.926);
\draw[rlAzulClaro, opacity=0.25] (1.326,1.170) -- (3.345,2.346);
\draw[rlAzulClaro, opacity=0.25] (1.326,1.170) -- (3.739,1.957);
\draw[rlAzulClaro, opacity=0.25] (1.326,1.170) -- (3.007,2.427);
\draw[rlAzulClaro, opacity=0.25] (1.326,1.170) -- (2.469,2.590);
\draw[rlAzulClaro, opacity=0.25] (1.326,1.170) -- (2.957,1.420);
\draw[rlAzulClaro, opacity=0.25] (1.326,1.170) -- (3.020,2.029);
\draw[rlAzulClaro, opacity=0.25] (1.326,1.170) -- (3.697,2.482);
\draw[rlAzulClaro, opacity=0.25] (1.326,1.170) -- (2.564,2.925);
\draw[rlAzulClaro, opacity=0.25] (1.326,1.170) -- (2.768,2.013);
\draw[rlAzulClaro, opacity=0.25] (1.326,1.170) -- (3.526,2.010);
\draw[rlAzulClaro, opacity=0.25] (1.326,1.170) -- (3.273,3.017);
\draw[rlAzulClaro, opacity=0.25] (1.326,1.170) -- (3.284,2.130);
\draw[rlAzulClaro, opacity=0.25] (1.326,1.170) -- (3.872,2.035);
\draw[rlAzulClaro, opacity=0.25] (1.326,1.170) -- (3.564,2.000);
\draw[rlAzulClaro, opacity=0.25] (1.326,1.170) -- (3.720,2.208);
\draw[rlAzulClaro, opacity=0.25] (1.326,1.170) -- (3.788,2.253);
\draw[rlAzulClaro, opacity=0.25] (1.326,1.170) -- (3.148,2.365);
\draw[rlAzulClaro, opacity=0.25] (1.326,1.170) -- (3.849,2.694);
\draw[rlAzulClaro, opacity=0.25] (1.326,1.170) -- (3.641,2.041);
\draw[rlAzulClaro, opacity=0.25] (1.326,1.170) -- (3.513,2.423);
\draw[rlAzulClaro, opacity=0.25] (1.326,1.170) -- (3.726,2.699);
\draw[rlAzulClaro, opacity=0.25] (1.326,1.170) -- (3.415,2.411);
\fill[rlAzulClaro, opacity=0.75] (3.285,2.400) circle (1.3pt);
\fill[rlAzulClaro, opacity=0.75] (3.179,2.253) circle (1.3pt);
\fill[rlAzulClaro, opacity=0.75] (3.315,1.861) circle (1.3pt);
\fill[rlAzulClaro, opacity=0.75] (3.316,2.651) circle (1.3pt);
\fill[rlAzulClaro, opacity=0.75] (2.923,2.269) circle (1.3pt);
\fill[rlAzulClaro, opacity=0.75] (3.362,2.408) circle (1.3pt);
\fill[rlAzulClaro, opacity=0.75] (3.515,1.746) circle (1.3pt);
\fill[rlAzulClaro, opacity=0.75] (3.108,1.926) circle (1.3pt);
\fill[rlAzulClaro, opacity=0.75] (3.345,2.346) circle (1.3pt);
\fill[rlAzulClaro, opacity=0.75] (3.739,1.957) circle (1.3pt);
\fill[rlAzulClaro, opacity=0.75] (3.007,2.427) circle (1.3pt);
\fill[rlAzulClaro, opacity=0.75] (2.469,2.590) circle (1.3pt);
\fill[rlAzulClaro, opacity=0.75] (2.957,1.420) circle (1.3pt);
\fill[rlAzulClaro, opacity=0.75] (3.020,2.029) circle (1.3pt);
\fill[rlAzulClaro, opacity=0.75] (3.697,2.482) circle (1.3pt);
\fill[rlAzulClaro, opacity=0.75] (2.564,2.925) circle (1.3pt);
\fill[rlAzulClaro, opacity=0.75] (2.768,2.013) circle (1.3pt);
\fill[rlAzulClaro, opacity=0.75] (3.526,2.010) circle (1.3pt);
\fill[rlAzulClaro, opacity=0.75] (3.273,3.017) circle (1.3pt);
\fill[rlAzulClaro, opacity=0.75] (3.284,2.130) circle (1.3pt);
\fill[rlAzulClaro, opacity=0.75] (3.872,2.035) circle (1.3pt);
\fill[rlAzulClaro, opacity=0.75] (3.564,2.000) circle (1.3pt);
\fill[rlAzulClaro, opacity=0.75] (3.720,2.208) circle (1.3pt);
\fill[rlAzulClaro, opacity=0.75] (3.788,2.253) circle (1.3pt);
\fill[rlAzulClaro, opacity=0.75] (3.148,2.365) circle (1.3pt);
\fill[rlAzulClaro, opacity=0.75] (3.849,2.694) circle (1.3pt);
\fill[rlAzulClaro, opacity=0.75] (3.641,2.041) circle (1.3pt);
\fill[rlAzulClaro, opacity=0.75] (3.513,2.423) circle (1.3pt);
\fill[rlAzulClaro, opacity=0.75] (3.726,2.699) circle (1.3pt);
\fill[rlAzulClaro, opacity=0.75] (3.415,2.411) circle (1.3pt);
\fill[rlAzulClaro, opacity=0.75] (3.317,1.731) circle (1.3pt);
\fill[rlAzulClaro, opacity=0.75] (3.564,1.861) circle (1.3pt);
\fill[rlAzulClaro, opacity=0.75] (3.209,2.580) circle (1.3pt);
\fill[rlAzulClaro, opacity=0.75] (3.617,2.022) circle (1.3pt);
\fill[rlAzulClaro, opacity=0.75] (3.555,2.029) circle (1.3pt);
\fill[rlAzulClaro, opacity=0.75] (3.653,2.153) circle (1.3pt);
\fill[rlAzulClaro, opacity=0.75] (3.505,2.570) circle (1.3pt);
\fill[rlAzulClaro, opacity=0.75] (3.190,1.477) circle (1.3pt);
\fill[rlAzulClaro, opacity=0.75] (3.541,2.045) circle (1.3pt);
\fill[rlAzulClaro, opacity=0.75] (3.232,1.903) circle (1.3pt);
\fill[rlAzulClaro, opacity=0.75] (3.707,1.753) circle (1.3pt);
\fill[rlAzulClaro, opacity=0.75] (3.703,2.197) circle (1.3pt);
\fill[rlAzulClaro, opacity=0.75] (2.783,2.641) circle (1.3pt);
\fill[rlAzulClaro, opacity=0.75] (3.084,2.403) circle (1.3pt);
\fill[rlAzulClaro, opacity=0.75] (3.509,1.762) circle (1.3pt);
\fill[rlAzulClaro, opacity=0.75] (3.556,2.523) circle (1.3pt);
\fill[rlAzulClaro, opacity=0.75] (2.725,2.456) circle (1.3pt);
\fill[rlAzulClaro, opacity=0.75] (3.612,2.565) circle (1.3pt);
\fill[rlAzulClaro, opacity=0.75] (3.063,2.131) circle (1.3pt);
\fill[rlAzulClaro, opacity=0.75] (3.156,2.312) circle (1.3pt);
\fill[rlAzulClaro, opacity=0.75] (3.302,1.828) circle (1.3pt);
\fill[rlAzulClaro, opacity=0.75] (3.453,2.497) circle (1.3pt);
\fill[rlAzulClaro, opacity=0.75] (3.755,1.821) circle (1.3pt);
\fill[rlAzulClaro, opacity=0.75] (2.980,1.748) circle (1.3pt);
\fill[rlAzulClaro, opacity=0.75] (3.441,2.426) circle (1.3pt);
\fill[rlAzulClaro, opacity=0.75] (3.313,1.700) circle (1.3pt);
\fill[rlAzulClaro, opacity=0.75] (3.463,1.663) circle (1.3pt);
\fill[rlAzulClaro, opacity=0.75] (3.534,2.582) circle (1.3pt);
\fill[rlAzulClaro, opacity=0.75] (2.239,2.664) circle (1.3pt);
\fill[rlAzulClaro, opacity=0.75] (3.196,2.038) circle (1.3pt);
\fill[rlAzulClaro, opacity=0.75] (3.520,1.863) circle (1.3pt);
\fill[rlAzulClaro, opacity=0.75] (2.984,1.697) circle (1.3pt);
\fill[rlAzulClaro, opacity=0.75] (3.246,1.907) circle (1.3pt);
\fill[rlAzulClaro, opacity=0.75] (3.186,2.161) circle (1.3pt);
\fill[rlAzulClaro, opacity=0.75] (2.877,2.278) circle (1.3pt);
\fill[rlAzulClaro, opacity=0.75] (3.053,2.677) circle (1.3pt);
\fill[rlAzulClaro, opacity=0.75] (3.230,2.824) circle (1.3pt);
\fill[rlAzulClaro, opacity=0.75] (3.255,2.252) circle (1.3pt);
\fill[rlAzulClaro, opacity=0.75] (3.384,2.058) circle (1.3pt);
\fill[rlAzulClaro, opacity=0.75] (2.989,2.215) circle (1.3pt);
\draw[-{Stealth[length=6pt]}, line width=1.6pt, rlAzul] (1.326,1.170) -- (3.354,2.184);
\fill[rlAzul] (1.326,1.170) circle (1.8pt);
\node[anchor=north west, font=\scriptsize\bfseries, text=rlAzulClaro] at (0.08,3.684) {Variância baixa};
\node[font=\scriptsize, text=rlAzul] at (2.574,1.365) {$\nabla J(\theta)$};
\node[anchor=south west, font=\tiny, text=rlLaranja] at (0.05,0.03) {região onde o passo piora o objetivo};
\draw[rlCinza!40] (0,0) rectangle (6.240,3.744);
\end{tikzpicture}

  \end{center}
  \vspace{-0.4em}
  {\small Cada ponto é a estimativa $\ghat$ de um lote; a seta é o gradiente verdadeiro.
  Os dois estimadores são \textbf{não-viesados}, com nuvens centradas em $\nabla J$.
  Mas, com variância alta, cerca de \alert{16\% dos passos} caem na região
  sombreada e pioram $J$; com variância baixa, praticamente nenhum.}

  \vspace{0.2em}
  {\scriptsize\color{rlCinza} Figura ilustrativa: 70 amostras gaussianas por painel, desvios-padrão 290 e 45
  em torno de um gradiente de norma $\approx 291$; fração na direção errada $\Phi(-\|\nabla J\|/\sigma)\approx 16\%$.}
\end{frame}

%% ============================================================
\begin{frame}{De onde vem a variância}
  \small
  \begin{tabular}{@{}p{0.22\textwidth}p{0.42\textwidth}p{0.29\textwidth}@{}}
    \toprule
    \textbf{Fonte de ruído} & \textbf{Mecanismo} & \textbf{Remédio (custo)} \\
    \midrule
    1. Aleatoriedade de $\tau$ & Transições e ações estocásticas variam a cada episódio & Inevitável sem modelo \\[0.3em]
    2. Crédito temporal & Recompensas anteriores a $t$ não foram causadas por $a_t$ & \azul{Trocar $R(\tau)$ por $G_t$ (grátis)} \\[0.3em]
    3. Escala do retorno & Se $G_t \approx 100$ sempre, todas as ações sobem juntas & \azul{Linha de base $b(s_t)$ (grátis)} \\[0.3em]
    4. Horizonte longo & $G_t$ soma muitas recompensas futuras, e o ruído acumula & {\color{rlLaranja}Crítico / TD (introduz viés)} \\
    \bottomrule
  \end{tabular}

  \vspace{1em}
  Esta aula aborda as linhas 3 e 4: linha de base e crítico.
\end{frame}

%% ============================================================
\section{Estimando a direção de subida}

\begin{frame}{A linha de base}
  \framesubtitle{Reduzir variância sem alterar o gradiente esperado}

  \begin{defbox}{gradiente de política com linha de base}
    Para qualquer função $b:\gS\to\R$ que dependa \textbf{apenas do estado}:
    \[
      \nabla_\theta J(\theta) = \E_{\tau\sim\pol_\theta}\Big[\sum_{t=0}^{\infty}\gamma^t\,
        \nabla_\theta\log\pol_\theta(a_t\mid s_t)\,\big(G_t - b(s_t)\big)\Big]
    \]
    Para toda escolha de $b(s_t)$ nessas condições, o estimador amostral continua \textbf{não-viesado}.
  \end{defbox}

  \begin{ideiabox}
    Aumentamos a log-probabilidade de $a_t$ na proporção em que o retorno obtido
    $G_t$ \textbf{superou a referência} $b(s_t)$ daquele estado.
  \end{ideiabox}
\end{frame}

%% ============================================================
\begin{frame}{Demonstração: a linha de base não introduz viés}
  \footnotesize
  \begin{align*}
    &\E_{\tau\sim\pol_\theta}\big[\nabla_\theta\log\pol_\theta(a_t\mid s_t)\,b(s_t)\big]
      && \text{\color{rlCinza}\itshape termo extra, para cada $t$ fixo}\\
    &\quad= \azul{\E_{s_{0:t},\,a_{0:t-1}}\Big[\,\E_{a_t,\,s_{t+1:\infty},\,a_{t+1:\infty}}}
       \big[\nabla_\theta\log\pol_\theta(a_t\mid s_t)\,b(s_t)\big]\azul{\Big]}
      && \text{\color{rlCinza}\itshape quebra da esperança}\\
    &\quad= \E_{s_{0:t},\,a_{0:t-1}}\Big[\,\azul{b(s_t)\;\E_{a_t\sim\pol_\theta(\cdot\mid s_t)}}
       \big[\nabla_\theta\log\pol_\theta(a_t\mid s_t)\big]\Big]
      && \text{\color{rlCinza}\itshape $b(s_t)$ sai da esperança interna}\\
    &\quad= \E_{s_{0:t},\,a_{0:t-1}}\Big[\,b(s_t)\azul{\sum_{a\in\gA}\pol_\theta(a\mid s_t)\,
       \frac{\nabla_\theta\pol_\theta(a\mid s_t)}{\pol_\theta(a\mid s_t)}}\Big]
      && \text{\color{rlCinza}\itshape razão de verossimilhança, ao contrário}\\
    &\quad= \E_{s_{0:t},\,a_{0:t-1}}\Big[\,b(s_t)\,\azul{\nabla_\theta\sum_{a\in\gA}\pol_\theta(a\mid s_t)}\Big]
       = \E\big[\,b(s_t)\,\azul{\nabla_\theta 1}\,\big] = \azul{0}
      && \text{\color{rlCinza}\itshape normalização: $\textstyle\sum_a\pol = 1$}\quad\blacksquare
  \end{align*}

  \begin{alertabox}[title={Onde a prova quebra se a linha de base depender da ação}]
    O passo decisivo é tirar $b(s_t)$ de dentro de $\sum_a$. Com $b(s_t,a_t)$ --- por exemplo
    $b = \Qpi(s,a)$, que zeraria o estimador! --- a fatoração é impossível e o estimador fica \textbf{viesado}.
  \end{alertabox}
\end{frame}

%% ============================================================
\begin{frame}{Por que centralizar o retorno ajuda}
  \framesubtitle{Em um ambiente com $+1$ por passo, todo retorno é alto}

  \begin{columns}[T]
    \begin{column}{0.48\textwidth}
      \begin{exbox}[title={Sem linha de base: o sinal bom some}]{}
        \small Uma ação medíocre num estado excelente recebe $G_t > 0$ alto e ainda
        assim é reforçada. A diferença entre ações boas e ruins é pequena em relação
        ao total de recompensas.
      \end{exbox}
    \end{column}
    \begin{column}{0.48\textwidth}
      \begin{teobox}{Com linha de base: medimos o mérito relativo}
        \small $G_t - b(s_t)$ mede quanto a ação ficou \textbf{acima ou abaixo do esperado}
        em $s_t$. Ações abaixo da média passam a ser desestimuladas.
      \end{teobox}
    \end{column}
  \end{columns}

  \vspace{0.8em}
  \begin{center}
    \textbf{Mesmo gradiente esperado, muito menos ruído por amostra.}
  \end{center}
\end{frame}

%% ============================================================
\begin{frame}{Linha de base ótima}
  \framesubtitle{Minimização da variância como mínimos quadrados ponderados}

  Para qualquer vetor aleatório $X$:
  \[
    \Var[X] = \E\big[\|X\|^2\big] - \big\|\E[X]\big\|^2
  \]
  Como $\E[X]$ \textbf{não depende de $b$} (prova anterior), minimizar a variância no estado $s$
  equivale a minimizar apenas o segundo momento:
  \[
    \arg\min_{b(s)}\;\E_{a\sim\pol_\theta(\cdot\mid s),\,G_t}\Big[\big\|\nabla_\theta\log\pol_\theta(a\mid s)\big\|^2\,
      \big(G_t-b(s)\big)^2\;\Big|\;s_t=s\Big]
  \]

  \begin{ideiabox}
    Um problema de mínimos quadrados em $b$, com pesos
    $w(a,s) = \big\|\nabla_\theta\log\pol_\theta(a\mid s)\big\|^2$.
  \end{ideiabox}
\end{frame}

%% ============================================================
\begin{frame}{A linha de base de variância mínima}
  \framesubtitle{Derivando a quadrática em relação ao escalar $b(s)$ e igualando a zero}

  \begin{teobox}{Resultado}
    \[
      b^*(s) = \frac{\E\big[\|\nabla_\theta\log\pol_\theta(a\mid s)\|^2\,G_t\mid s_t=s\big]}
                    {\E\big[\|\nabla_\theta\log\pol_\theta(a\mid s)\|^2\mid s_t=s\big]}
      \;\approx\;\E_{\pol_\theta}[G_t\mid s_t=s] = \Vpi(s)
    \]
  \end{teobox}

  \begin{itemize}
    \item $b^*(s)$ é a média dos retornos \textbf{ponderada} pelo quadrado da norma do escore
          $\|\nabla_\theta\log\pol_\theta\|^2$.
    \item Ignorando essa ponderação (o padrão em implementações), obtém-se exatamente $\Vpi(s)$.
  \end{itemize}

  \begin{center}
    \textbf{Linha de base prática e crítico de valor coincidem.}
  \end{center}
\end{frame}

%% ============================================================
\begin{frame}{Exercício resolvido: bandido softmax de três ações}
  \framesubtitle{Um cálculo fechado: quanto se ganha exatamente com a linha de base?}

  Um único estado, três ações, política softmax e recompensas determinísticas:
  \[
    \theta=(0,0,0)\;\Rightarrow\;\pol = \big(\tfrac13,\tfrac13,\tfrac13\big),\qquad r = (98,\,101,\,104)
  \]
  O escore da softmax tem a mesma norma para toda ação $a$:
  \[
    \nabla_\theta\log\pol(a) = e_a - \pol\;\;\Rightarrow\;\;
    \|\nabla_\theta\log\pol(a)\|^2 = \big(1-\tfrac13\big)^2 + 2\big(0-\tfrac13\big)^2 = \tfrac23
  \]
  Pesos idênticos, logo a linha de base ótima é a média simples, $b^* = 101$. Estimador de uma amostra:
  \[
    \ghat = (e_a-\pol)\,\big(r(a)-b\big),\qquad
    \mathrm{tr}\,\mathrm{Cov}[\ghat] = \E\|\ghat\|^2 - \|\E[\ghat]\|^2
  \]
\end{frame}

%% ============================================================
\begin{frame}{Resultado: a variância cai 3400$\times$}
  \framesubtitle{Mesmo gradiente esperado, sem introduzir viés}

  \begin{center}
    \begin{tabular}{@{}lccr@{}}
      \toprule
      \textbf{Configuração} & $\E[\ghat]$ & $\E\|\ghat\|^2$ & \textbf{Variância total} \\
      \midrule
      Sem linha de base $(b=0)$   & $(-1,\,0,\,+1)$ & 6804,67 & {\color{rlLaranja}6802,67} \\
      Com linha de base $(b=101)$ & $(-1,\,0,\,+1)$ & 4,00    & \azul{2,00} \\
      \bottomrule
    \end{tabular}
  \end{center}

  \vspace{0.5em}
  Sem linha de base, seriam necessárias cerca de \textbf{3400$\times$ mais trajetórias}
  para o mesmo erro-padrão $(\propto 1/\sqrt{m})$.

  \begin{ideiabox}[title={Por quê?}]
    Sem linha de base, cada amostra carrega o nível comum das recompensas ($\approx 101$).
    Com a linha de base em 101, sobra só a diferença entre as ações ($\pm 3$). Como a variância
    cresce com o quadrado dessa escala, o ganho é da ordem de (nível médio / dispersão)$^2$.
  \end{ideiabox}
\end{frame}

%% ============================================================
\begin{frame}{Do retorno amostral ao crítico}
  \framesubtitle{Trocar a soma longa $G_t$ por bootstrapping e função vantagem}

  Mesmo em $G_t - \Vpi(s_t)$, o alvo $G_t$ soma recompensas aleatórias até o fim do episódio
  (Monte Carlo), com variância que cresce com o horizonte. A saída replica a passagem de MC
  para TD (Aula 3): \emph{bootstrapping}.

  \begin{defbox}{métodos ator-crítico e função vantagem}
    Mantêm representações paramétricas explícitas do \textbf{ator} $\pol_\theta$ e do
    \textbf{crítico} (valor $V$ ou $Q$ com parâmetros $w$), atualizadas em paralelo:
    \[
      \nabla_\theta J(\theta) \approx \E_{\tau\sim\pol_\theta}\Big[\sum_{t=0}^{\infty}\gamma^t\,
        \nabla_\theta\log\pol_\theta(a_t\mid s_t)\,\Ahat^{\pol}(s_t,a_t)\Big],
      \qquad
      A^{\pol}(s,a) \defeq \Qpi(s,a) - \Vpi(s)
    \]
  \end{defbox}
\end{frame}

%% ============================================================
\begin{frame}{A vantagem é centrada por construção}
  \framesubtitle{Por que $A^{\pol}$ é o sinal de reforço ``certo''}

  \begin{teobox}{Propriedade de centralização}
    Para todo $s\in\gS$, como $\Vpi(s) = \E_{a\sim\pol}\big[\Qpi(s,a)\big]$:
    \[
      \E_{a\sim\pol(\cdot\mid s)}\big[A^{\pol}(s,a)\big] = \Vpi(s) - \Vpi(s) = 0
    \]
  \end{teobox}

  \begin{columns}[T]
    \begin{column}{0.48\textwidth}
      \begin{ideiabox}[title={Sinal de média zero}]
        Na média ponderada por $\pol(\cdot\mid s)$, ações acima da média recebem peso
        positivo e ações abaixo da média, peso negativo.
      \end{ideiabox}
    \end{column}
    \begin{column}{0.48\textwidth}
      \begin{ideiabox}[title={Ligação com a linha de base}]
        $A^{\pol} = \Qpi - \Vpi$ é ``retorno esperado menos a linha de base prática'':
        ator-crítico = linha de base + \emph{bootstrapping}.
      \end{ideiabox}
    \end{column}
  \end{columns}
\end{frame}

%% ============================================================
\begin{frame}{O espectro de estimadores: de Monte Carlo a TD}
  \footnotesize\setlength{\tabcolsep}{4pt}
  \begin{tabular}{@{}llll@{}}
    \toprule
    \textbf{Peso no estimador do gradiente} & \textbf{Nome} & \textbf{Viés} & \textbf{Variância} \\
    \midrule
    $R(\tau)$ & Retorno total da trajetória & Nenhum & Altíssima \\
    $G_t$ & Retorno futuro (causal) & Nenhum & Muito alta \\
    $G_t - b(s_t)$ & Retorno com linha de base analítica & Nenhum & Alta \\
    $G_t - V_w(s_t)$ & Vantagem de Monte Carlo & Nenhum$^\dagger$ & Média \\
    $Q_w(s_t,a_t) - V_w(s_t)$ & Vantagem estimada dupla & Do crítico $Q$ & Baixa \\
    $\delta_t = r_t + \gamma V_w(s_{t+1}) - V_w(s_t)$ & Erro de TD de 1 passo & Do crítico $V$ & Muito baixa \\
    \rowcolor{rlAzulClaro!8}
    $\Ahat_t^{\mathrm{GAE}(\gamma,\lambda)} = \sum_{l\ge0}(\gamma\lambda)^l\,\delta_{t+l}$ & Vantagem generalizada (Aula 7) & Ajustável $(\lambda)$ & Ajustável $(\lambda)$ \\
    \bottomrule
  \end{tabular}

  \vspace{0.8em}
  {\footnotesize $^\dagger$ Subtrair $V(s_t;w)$ de $G_t$ não introduz viés no gradiente esperado (desde que
  $w$ não dependa da ação $a_t$), embora erros no crítico degradem a redução de variância.}
\end{frame}

%% ============================================================
\begin{frame}{A variância de cada estimador ao longo do treino}
  \framesubtitle{Cadeia de 8 estados com recompensa ruidosa; todos medidos no mesmo lote}

  \begin{center}
    % Gerado por figs.py a partir da simulação descrita no anexo (média móvel de 5 iterações).
\providecommand{\figVarW}{0.66\textwidth}\providecommand{\figVarH}{0.5\textheight}
\begin{tikzpicture}
\begin{axis}[width=\figVarW, height=\figVarH, ymode=log,
  xmin=0, xmax=146, ymin=60, ymax=60000, xlabel={iteração de treino}, ylabel={variância do estimador},
  tick label style={font=\scriptsize}, label style={font=\scriptsize}, axis lines=left,
  ymajorgrids, grid style={rlCinza!20}, legend pos=outer north east, legend style={font=\scriptsize, draw=none, fill=none},
  legend cell align=left, every axis plot/.append style={line width=1.4pt, no marks}]
\addplot[rlLaranja] coordinates {(0,27858) (4,28946) (8,27533) (12,28728) (16,30075) (20,24702) (24,26997) (28,25078) (32,28701) (36,28448) (40,28790) (44,24869) (48,28547) (52,28117) (56,27861) (60,26534) (64,26224) (68,29416) (72,24915) (76,26646) (80,26412) (84,26305) (88,30993) (92,26904) (96,28859) (100,27413) (104,27191) (108,26759) (112,25432) (116,26474) (120,30453) (124,31224) (128,29334) (132,26509) (136,29296) (140,27754) (144,25973) (145,26860)};
\addlegendentry{$R(\tau)$}
\addplot[rlAmbar] coordinates {(0,16704) (4,16736) (8,16802) (12,17568) (16,17375) (20,13964) (24,16038) (28,14040) (32,17350) (36,17341) (40,17269) (44,15435) (48,18374) (52,17844) (56,17354) (60,16667) (64,16703) (68,16907) (72,15117) (76,16420) (80,16061) (84,15556) (88,18400) (92,16159) (96,17828) (100,16869) (104,17048) (108,16878) (112,16542) (116,16949) (120,17696) (124,18742) (128,17729) (132,17427) (136,17515) (140,17840) (144,15772) (145,16010)};
\addlegendentry{$G_t$}
\addplot[rlAzulClaro!55] coordinates {(0,7228) (4,4889) (8,4502) (12,4802) (16,5089) (20,4488) (24,4375) (28,4039) (32,4978) (36,4842) (40,4388) (44,4063) (48,4566) (52,4783) (56,4920) (60,4661) (64,4427) (68,4805) (72,4659) (76,5341) (80,5009) (84,4547) (88,5167) (92,4450) (96,4923) (100,4930) (104,4446) (108,4766) (112,4520) (116,5247) (120,4784) (124,4765) (128,5217) (132,4891) (136,4769) (140,4482) (144,4456) (145,4585)};
\addlegendentry{$G_t - V(s_t)$}
\addplot[rlAzulClaro] coordinates {(0,3219) (4,2178) (8,1940) (12,1915) (16,1900) (20,1742) (24,1878) (28,1734) (32,1821) (36,1795) (40,1858) (44,1644) (48,1670) (52,1648) (56,1656) (60,1563) (64,1573) (68,1823) (72,1662) (76,1811) (80,1594) (84,1830) (88,2132) (92,1855) (96,1855) (100,1721) (104,1655) (108,1622) (112,1616) (116,1626) (120,1797) (124,1881) (128,1692) (132,1667) (136,1724) (140,1636) (144,1615) (145,1660)};
\addlegendentry{GAE, $\lambda=0{,}95$}
\addplot[rlAzul] coordinates {(0,457) (4,293) (8,192) (12,161) (16,136) (20,120) (24,144) (28,139) (32,114) (36,136) (40,162) (44,131) (48,158) (52,133) (56,130) (60,145) (64,189) (68,189) (72,150) (76,169) (80,177) (84,196) (88,210) (92,172) (96,166) (100,123) (104,121) (108,144) (112,159) (116,142) (120,164) (124,174) (128,148) (132,163) (136,158) (140,162) (144,135) (145,129)};
\addlegendentry{TD, $\delta_t$}
\end{axis}
\end{tikzpicture}

  \end{center}
  \vspace{-0.6em}
  {\small Escala log: cada linha de grade é 10$\times$ mais variância. As curvas com crítico caem
  no início, enquanto o crítico aprende; a ordem é a da tabela anterior.}
  {\scriptsize\color{rlCinza} Problema e código: anexo, último slide.}
\end{frame}

%% ============================================================
\begin{frame}{O preço explícito do erro de TD}
  \framesubtitle{Onde o bootstrapping troca variância por viés}

  \begin{alertabox}[title={O erro de TD só é não-viesado com o crítico exato}]
    \[
      \E[\delta_t\mid s_t,a_t] = A^{\pol}(s_t,a_t)\quad\text{apenas quando } V_w = \Vpi
    \]
    Com um crítico imperfeito, surge o viés:
    \[
      \E[\delta_t\mid s_t,a_t] - A^{\pol}(s_t,a_t) = \gamma\,\E\big[V_w(s_{t+1}) - \Vpi(s_{t+1})\mid s_t,a_t\big]
        - \big(V_w(s_t)-\Vpi(s_t)\big)
    \]
  \end{alertabox}

  \begin{ideiabox}
    Troca deliberada: aceitamos um viés controlado pela qualidade do crítico em troca de
    variância mínima. A GAE (Aula 7) torna essa troca um botão contínuo, $\lambda$.
  \end{ideiabox}
\end{frame}

%% ============================================================
\section{Geometria do passo e eficiência amostral}

\begin{frame}{As duas barreiras do gradiente \emph{on-policy}}
  \framesubtitle{Limitações que persistem mesmo com a variância já controlada}

  \begin{columns}[T]
    \begin{column}{0.48\textwidth}
      \begin{teobox}{1 · Ineficiência amostral: os dados expiram a cada passo}
        \small O gradiente é uma esperança sob $\pol_\theta$. Assim que fazemos
        $\theta\leftarrow\theta+\alpha\,\ghat$, os dados coletados descrevem uma política que
        \textbf{não existe mais}, e o lote inteiro é descartado.\\[0.3em]
        \textbf{Meta do PPO:} quantos passos podemos dar no mesmo lote antes que os dados percam validade?
      \end{teobox}
    \end{column}
    \begin{column}{0.48\textwidth}
      \begin{teobox}{2 · Parâmetros vs. políticas: a régua euclidiana engana}
        \small No caso finito, as políticas vivem em
        \[
          \Pi = \Big\{\pol\in\R^{|\gS|\times|\gA|} : \textstyle\sum_a \pol(a\mid s) = 1,\ \pol\ge 0\Big\}
        \]
        Damos passos em $\theta\in\Theta$, mas $J(\pol_\theta)$ depende do deslocamento real em
        $\Pi$, e a mesma régua $\|\Delta\theta\|$ não serve nos dois espaços.
      \end{teobox}
    \end{column}
  \end{columns}
\end{frame}

%% ============================================================
\begin{frame}{Consequência prática: o colapso de desempenho}
  \begin{center}
    {\Large\bfseries\color{rlAzul} Em RL, um passo grande demais destrói a política.}
  \end{center}
  \vspace{0.6em}
  No aprendizado supervisionado os dados são fixos e o próximo lote corrige um passo ruim.
  Em RL, a política ruim passa a coletar \alert{apenas trajetórias ruins}, das quais o agente
  nunca mais aprende a escapar.

  \vspace{0.6em}
  E um passo fixo minúsculo torna o treino inviável. Precisamos medir o passo no espaço certo.
\end{frame}

%% ============================================================
\begin{frame}{Distância entre políticas: divergência KL}
  \framesubtitle{O mesmo $\Delta\theta = 2$ pode ser catastrófico ou quase nulo}

  \begin{defbox}{Kullback--Leibler entre políticas num estado}
    \[
      \KL(\pol'\,\|\,\pol)[s] \defeq \sum_{a\in\gA}\pol'(a\mid s)\,\log\frac{\pol'(a\mid s)}{\pol(a\mid s)}
    \]
    \[
      \KL(P\,\|\,Q)\ge 0,\qquad \KL(P\,\|\,P) = 0,\qquad \KL(P\,\|\,Q)\neq\KL(Q\,\|\,P)\ \text{ em geral}
    \]
  \end{defbox}

  \begin{ideiabox}
    Ao contrário de $\|\Delta\theta\|$, a KL mede o quanto o \textbf{comportamento} mudou:
    probabilidade de ações, estado a estado. O $[s]$ indica o estado em que a KL é calculada.
  \end{ideiabox}
\end{frame}

%% ============================================================
\begin{frame}{Exercício resolvido: sigmoide com um único parâmetro escalar}
  \framesubtitle{Dois passos idênticos $\Delta\theta = 2$, em regiões diferentes}

  \[
    \pol_\theta(a{=}1) = \sigma(\theta) = \frac{1}{1+e^{-\theta}},\qquad \pol_\theta(a{=}2) = 1-\sigma(\theta)
  \]
  \begin{center}\small
    \begin{tabular}{@{}lrrrr@{}}
      \toprule
      \textbf{Passo} & $\pol_{\text{ini}}(a{=}1)$ & $\pol_{\text{fim}}(a{=}1)$ & $|\Delta\pol|$ & $\KL(\pol_{\text{fim}}\,\|\,\pol_{\text{ini}})$ \\
      \midrule
      Região linear: $\theta: -1\to+1$ & 0,2689 & 0,7311 & \azul{0,4621} & \azul{0,4621} \\
      Região saturada: $\theta: 4\to6$ & 0,9820 & 0,9975 & {\color{rlLaranja}0,0155} & {\color{rlLaranja}0,0107} \\
      \bottomrule
    \end{tabular}
  \end{center}

  O mesmo $\Delta\theta$ muda o comportamento por fatores de \textbf{30$\times$ em probabilidade}
  e \textbf{43$\times$ em KL}. Nenhuma taxa de aprendizado constante serve para toda a trajetória de treino.
\end{frame}

%% ============================================================
\begin{frame}{Geometria local: visitação descontada}
  \framesubtitle{Restringir a KL equivale a usar a régua certa em cada ponto $\theta$}

  \begin{defbox}{distribuição descontada de visitação de estados}
    \[
      d^{\pol}(s) \defeq (1-\gamma)\sum_{t=0}^{\infty}\gamma^t\,\Prob(s_t = s\mid\pol),\qquad
      \sum_{s\in\gS} d^{\pol}(s) = (1-\gamma)\sum_{t=0}^{\infty}\gamma^t = 1
    \]
  \end{defbox}

  Para deslocamentos pequenos $\Delta\theta$, a expansão de Taylor de 2ª ordem da KL média nos
  estados $s\sim d^{\pol_\theta}$ dá:
  \[
    \KLbar(\theta+\Delta\theta\,\|\,\theta) \defeq \E_{s\sim d^{\pol_\theta}}\big[\KL(\pol_{\theta+\Delta\theta}\,\|\,\pol_\theta)[s]\big]
      \;\approx\;\tfrac12\,\Delta\theta^\top F(\theta)\,\Delta\theta
  \]
  em que a curvatura é dada pela matriz de informação de Fisher:
  \[
    F(\theta) = \E_{s\sim d^{\pol_\theta},\,a\sim\pol_\theta}\Big[\nabla_\theta\log\pol_\theta(a\mid s)\,
      \nabla_\theta\log\pol_\theta(a\mid s)^\top\Big]
  \]
\end{frame}

%% ============================================================
\begin{frame}{A matriz de informação de Fisher}
  \framesubtitle{A métrica local do espaço de políticas}

  \begin{defbox}{}
    \[
      F(\theta) = \E_{s\sim d^{\pol_\theta},\,a\sim\pol_\theta}\Big[\nabla_\theta\log\pol_\theta(a\mid s)\,
        \nabla_\theta\log\pol_\theta(a\mid s)^\top\Big]
    \]
    Pequenos passos: $\KLbar\approx\tfrac12\,\Delta\theta^\top F(\theta)\,\Delta\theta$.
    Fisher é a régua de $\Pi$ vista de $\Theta$.
  \end{defbox}

  \begin{columns}[T]
    \begin{column}{0.48\textwidth}
      \begin{teobox}{TRPO (Schulman et al., 2015): segunda ordem explícita}
        \small Usa o gradiente natural $F(\theta)^{-1}g$ sob uma restrição quadrática de KL.
      \end{teobox}
    \end{column}
    \begin{column}{0.48\textwidth}
      \begin{exbox}[title={PPO (Schulman et al., 2017): apenas primeira ordem}]{}
        \small Impõe proximidade em $\Pi$ sem construir nem inverter $F(\theta)$.
      \end{exbox}
    \end{column}
  \end{columns}
\end{frame}

%% ============================================================
\begin{frame}{Lema da diferença de desempenho}
  \framesubtitle{A identidade fundamental de Trust Region Policy Optimization (TRPO) e Proximal Policy Optimization (PPO)}

  \begin{teobox}{Lema (Kakade \& Langford, 2002)}
    Para quaisquer duas políticas $\pol$ (antiga) e $\pol'$ (nova):
    \[
      J(\pol') - J(\pol) = \E_{\tau\sim\pol'}\Big[\sum_{t=0}^{\infty}\gamma^t A^{\pol}(s_t,a_t)\Big]
        = \frac{1}{1-\gamma}\,\E_{s\sim d^{\pol'},\,a\sim\pol'}\big[A^{\pol}(s,a)\big]
    \]
  \end{teobox}

  \begin{ideiabox}
    A nova política melhora exatamente na medida em que escolhe, \textbf{nos estados que ela própria visita},
    ações com vantagem positiva segundo a política antiga.
  \end{ideiabox}
\end{frame}

%% ============================================================
\begin{frame}{Demonstração do lema}
  Escrevendo a vantagem como resíduo de Bellman:
  \[
    A^{\pol}(s_t,a_t) = \E_{s_{t+1}}\big[r_t + \gamma\Vpi(s_{t+1}) - \Vpi(s_t)\big]
  \]
  Somando ao longo de uma trajetória da política nova:
  \begin{align*}
    \E_{\tau\sim\pol'}\Big[\sum_{t=0}^{\infty}\gamma^t A^{\pol}(s_t,a_t)\Big]
      &= \E_{\tau\sim\pol'}\Big[\sum_{t=0}^{\infty}\gamma^t r_t\Big]
       + \E_{\tau\sim\pol'}\Big[\sum_{t=0}^{\infty}\big(\gamma^{t+1}\Vpi(s_{t+1}) - \gamma^t\Vpi(s_t)\big)\Big]\\
      &= J(\pol') + \E_{s_0}\big[-\Vpi(s_0)\big] = J(\pol') - J(\pol)
  \end{align*}
  A segunda soma é \textbf{telescópica}: sobra apenas o termo de $t=0$, e
  $\E[\Vpi(s_0)] = J(\pol)$. Pela definição de $d^{\pol'}$, segue a forma compacta. $\blacksquare$
\end{frame}

%% ============================================================
\begin{frame}{Tornando o ganho computável}
  \framesubtitle{Eliminando a dependência de trajetórias da política nova $\pol'$}

  No lema exato, $A^{\pol}$ já usa a régua antiga, mas ações e estados ainda vêm de $\pol'$.
  Fazemos duas operações:

  \medskip
  {\small\bfseries\azul{1 · Trocar $a\sim\pol'$ por $a\sim\pol$} (exato, por amostragem por importância em cada passo)}
  \[
    J(\pol') - J(\pol) = \frac{1}{1-\gamma}\,\E_{s\sim d^{\pol'},\,a\sim\pol}\Big[\frac{\pol'(a\mid s)}{\pol(a\mid s)}\,A^{\pol}(s,a)\Big]
  \]
  {\small\bfseries{\color{rlLaranja}2 · Trocar estados $d^{\pol'}$ por $d^{\pol}$} (aproximação local de 1ª ordem)}
  \[
    L_{\pol}(\pol') \defeq \frac{1}{1-\gamma}\,\E_{s\sim d^{\pol},\,a\sim\pol}\Big[\frac{\pol'(a\mid s)}{\pol(a\mid s)}\,A^{\pol}(s,a)\Big]
      = \E_{\tau\sim\pol}\Big[\sum_{t=0}^{\infty}\gamma^t\,\frac{\pol'(a_t\mid s_t)}{\pol(a_t\mid s_t)}\,A^{\pol}(s_t,a_t)\Big]
  \]
\end{frame}

%% ============================================================
\begin{frame}{Por que os pesos não explodem}
  \framesubtitle{E o que ganhamos com isso}

  \begin{columns}[T]
    \begin{column}{0.48\textwidth}
      \begin{exbox}[title={Importância de trajetórias inteiras}]{}
        \small O peso é um produto, que explode ou zera exponencialmente com $t$:
        \[
          \prod_{t'=0}^{t}\frac{\pol'(a_{t'}\mid s_{t'})}{\pol(a_{t'}\mid s_{t'})}
        \]
      \end{exbox}
    \end{column}
    \begin{column}{0.48\textwidth}
      \begin{teobox}{No objetivo substituto}
        \small Aparece \textbf{uma única razão por passo}:
        \[
          \frac{\pol'(a_t\mid s_t)}{\pol(a_t\mid s_t)}
        \]
      \end{teobox}
    \end{column}
  \end{columns}

  \begin{ideiabox}
    Tudo é estimável com o lote antigo $\tau\sim\pol$, o que habilita
    \textbf{vários passos de gradiente por coleta}.
  \end{ideiabox}
\end{frame}

%% ============================================================
\begin{frame}{Quando $L_{\pol}(\pol')$ é confiável?}
  \framesubtitle{Aproximação local exige um mecanismo explícito de região de validade}

  Em $\pol'=\pol$, o substituto toca o objetivo verdadeiro com a inclinação exata:
  \[
    L_{\pol}(\pol) = 0,\qquad \nabla_{\theta'}L_{\pol_\theta}(\pol_{\theta'})\big|_{\theta'=\theta} = \nabla_\theta J(\pol_\theta)
  \]

  \begin{teobox}{Teorema (Achiam, Held, Tamar \& Abbeel, 2017)}
    Existe $C>0$, dependente de $\gamma$ e de $\max_{s,a}|A^{\pol}(s,a)|$, tal que, para quaisquer $\pol$ e $\pol'$:
    \[
      \Big|\,J(\pol') - \big(J(\pol) + L_{\pol}(\pol')\big)\Big| \;\le\;
        C\,\sqrt{\E_{s\sim d^{\pol}}\big[\KL(\pol'\,\|\,\pol)[s]\big]}
    \]
  \end{teobox}
\end{frame}

%% ============================================================
\begin{frame}{Maximizar um limite inferior}
  \framesubtitle{A licença para reutilizar o lote é condicional}

  \begin{teobox}{Consequência do teorema}
    Maximizar o limitante inferior garante melhoria real em $J(\pol)$:
    \[
      J(\pol') \;\ge\; J(\pol) + L_{\pol}(\pol') - C\sqrt{\KLbar(\pol'\,\|\,\pol)}
    \]
  \end{teobox}

  \begin{alertabox}[title={Mas atenção}]
    O teorema vale globalmente, mas só é informativo quando a KL média é pequena.
    Maximizar $L_{\pol}(\pol')$ \textbf{sem restrição} empurra a razão
    $\pol'(a\mid s)/\pol(a\mid s)\gg 1$ exatamente para longe de $\pol$, otimizando o modelo
    fora de sua região de validade.
  \end{alertabox}
\end{frame}

%% ============================================================
\begin{frame}{PPO, variante 1: penalidade KL adaptativa}
  \framesubtitle{Convertendo a restrição de proximidade em regularização autoajustada}

  Transforma-se o problema restrito numa otimização sem restrições, penalizada pela KL empírica:
  \[
    \theta_{k+1} = \arg\max_\theta\Big(L_{\theta_k}(\theta) - \beta_k\,\KLbar(\theta_k,\theta)\Big)
  \]
  \[
    \KLbar(\theta_k,\theta) \defeq \E_{s\sim d^{\pol_k}}\Big[\KL\big(\pol_{\theta_k}(\cdot\mid s)\,\|\,\pol_\theta(\cdot\mid s)\big)\Big]
  \]
  {\small\color{rlCinza} Nota: tomar $\pol_{\theta_k}$ no 1º argumento facilita o cálculo com dados de
  $\pol_{\theta_k}$ e preserva a mesma curvatura local $\tfrac12\,\Delta\theta^\top F(\theta_k)\,\Delta\theta$.}
\end{frame}

%% ============================================================
\begin{frame}{Algoritmo: PPO com penalidade KL adaptativa}
  \begin{algbox}{PPO com penalidade KL adaptativa}
    \small
    \textbf{Entrada:} $\theta_0$, penalidade inicial $\beta_0$, divergência alvo $\delta$, épocas $K$.\\
    \textbf{para} $k = 0, 1, 2, \dots$ \textbf{faça}
    \begin{enumerate}\setlength{\itemsep}{0pt}
      \item Colete um lote de trajetórias $\mathcal{D}_k$ executando a política atual $\pol_{\theta_k}$.
      \item Estime as vantagens $\Ahat_t$ usando um crítico de valor.
      \item Otimize $L_{\theta_k}(\theta) - \beta_k\,\KLbar(\theta_k,\theta)$ por \textbf{$K$ épocas de SGD/Adam} em minilotes.
      \item Adapte o peso: se $\KLbar\ge 1{,}5\,\delta$, $\beta_{k+1}\leftarrow 2\beta_k$ (afastou-se demais);
            se $\KLbar\le\delta/1{,}5$, $\beta_{k+1}\leftarrow\beta_k/2$ (cauteloso demais); senão, mantenha.
    \end{enumerate}
  \end{algbox}
\end{frame}

%% ============================================================
\begin{frame}{PPO, variante 2: objetivo recortado}
  \framesubtitle{Removendo o incentivo para se afastar de $r_t(\theta) = 1$, sem calcular KL}

  \begin{defbox}{razão de probabilidades e objetivo recortado}
    \[
      r_t(\theta) \defeq \frac{\pol_\theta(a_t\mid s_t)}{\pol_{\theta_k}(a_t\mid s_t)}\qquad(\text{logo } r_t(\theta_k) = 1)
    \]
    Com hiperparâmetro $\epsilon>0$ (tipicamente $\epsilon = 0{,}2$):
    \[
      L^{\mathrm{CLIP}}_{\theta_k}(\theta) = \E_{\tau\sim\pol_{\theta_k}}\Big[\sum_{t=0}^{T-1}
        \min\Big(r_t(\theta)\,\Ahat_t,\ \clipop\big(r_t(\theta),1-\epsilon,1+\epsilon\big)\,\Ahat_t\Big)\Big]
    \]
  \end{defbox}
\end{frame}

%% ============================================================
\begin{frame}{Leitura gráfica do recorte}
  \framesubtitle{Um termo de $L^{\mathrm{CLIP}}$ em função de $r_t$, para cada sinal da vantagem}

  \begin{columns}[T]
    \begin{column}{0.48\textwidth}
      \centering{\small\bfseries\azul{Caso $\Ahat_t>0$ (ação boa)}}\\[0.3em]
      \begin{tikzpicture}[x=2.6cm, y=1.5cm]
        \draw[-{Stealth}] (0,0) -- (2.15,0) node[below left, font=\scriptsize] {$r_t(\theta)$};
        \draw[-{Stealth}] (0,0) -- (0,1.75) node[right, font=\scriptsize, rlCinza] {objetivo};
        \foreach \x/\l in {0.8/1-\epsilon, 1/1, 1.2/1+\epsilon}
          \draw[rlCinza!60, densely dotted] (\x,0) -- (\x,\x) node[pos=0, below, font=\tiny, text=black] {$\l$};
        \draw[rlCinza!70, dashed, line width=1pt] (1.2,1.2) -- (1.6,1.6) node[right, font=\scriptsize, text=rlCinza] {$r_t\Ahat_t$};
        \draw[rlAzulClaro, line width=2.2pt] (0,0) -- (1.2,1.2) -- (2.05,1.2);
        \fill[rlAzul] (1,1) circle (2pt);
        \node[font=\scriptsize, text=rlAzulClaro, above] at (1.65,1.2) {$\nabla_\theta = 0$ (teto)};
      \end{tikzpicture}
    \end{column}
    \begin{column}{0.48\textwidth}
      \centering{\small\bfseries{\color{rlLaranja}Caso $\Ahat_t<0$ (ação ruim)}}\\[0.3em]
      \begin{tikzpicture}[x=2.6cm, y=1.5cm]
        \draw[-{Stealth}] (0,0) -- (2.15,0) node[above left, font=\scriptsize] {$r_t(\theta)$};
        \draw (0,0.1) -- (0,-1.75) node[right, font=\scriptsize, rlCinza] {objetivo};
        \foreach \x/\l in {0.8/1-\epsilon, 1/1, 1.2/1+\epsilon}
          \draw[rlCinza!60, densely dotted] (\x,0) -- (\x,-\x) node[pos=0, above, font=\tiny, text=black] {$\l$};
        \draw[rlCinza!70, dashed, line width=1pt] (0,0) -- (0.8,-0.8);
        \node[font=\scriptsize, text=rlCinza] at (0.25,-0.5) {$r_t\Ahat_t$};
        \draw[rlLaranja, line width=2.2pt] (0,-0.8) -- (0.8,-0.8) -- (2.0,-2.0+0.4);
        \fill[rlAzul] (1,-1) circle (2pt);
        \node[font=\scriptsize, text=rlLaranja, below] at (0.4,-0.8) {$\nabla_\theta = 0$ (piso)};
      \end{tikzpicture}
    \end{column}
  \end{columns}

  \vspace{0.3em}
  {\small\textbf{Eixo vertical:} o valor do termo de uma amostra,
  $\min\big(r_t\Ahat_t,\ \clipop(r_t,1-\epsilon,1+\epsilon)\,\Ahat_t\big)$ (o eixo horizontal marca o zero).
  Tracejado: o termo sem recorte; o ponto marca $r_t = 1$.}
\end{frame}

%% ============================================================
\begin{frame}{A assimetria deliberada do recorte}
  \framesubtitle{Por que o $\min$ atua como limitante inferior pessimista}

  \begin{center}\footnotesize\setlength{\tabcolsep}{4pt}
    \begin{tabular}{@{}p{0.16\textwidth}p{0.25\textwidth}p{0.19\textwidth}p{0.3\textwidth}@{}}
      \toprule
      \textbf{Sinal da vantagem} & $r_t < 1-\epsilon$ & $1-\epsilon\le r_t\le 1+\epsilon$ & $r_t > 1+\epsilon$ \\
      \midrule
      $\Ahat_t>0$ (ação boa) & $r_t\Ahat_t$ --- ativo: recupera & $r_t\Ahat_t$ --- ativo &
        {\color{rlLaranja}$(1+\epsilon)\Ahat_t$ --- plano: trava a subida} \\[0.3em]
      $\Ahat_t<0$ (ação ruim) & {\color{rlLaranja}$(1-\epsilon)\Ahat_t$ --- plano: trava a queda} &
        $r_t\Ahat_t$ --- ativo & \azul{$r_t\Ahat_t$ --- ativo: corrige o excesso} \\
      \bottomrule
    \end{tabular}
  \end{center}

  \begin{ideiabox}
    O recorte só desliga o gradiente quando a mudança já foi \textbf{na direção que a vantagem pedia}.
    Movimentos na direção errada continuam sempre corrigíveis.
  \end{ideiabox}
\end{frame}

%% ============================================================
\begin{frame}{O que o recorte faz, e o que não faz}
  \framesubtitle{Duas leituras para guardar}

  \begin{columns}[T]
    \begin{column}{0.48\textwidth}
      \begin{teobox}{Freia alteração, nunca o arrependimento}
        \small Para $\Ahat_t>0$, o recorte desliga o gradiente acima de $1+\epsilon$.\\[0.3em]
        Caso uma ação ruim $(\Ahat_t<0)$ tenha sido acidentalmente empurrada para $r_t>1+\epsilon$
        por minilotes anteriores, o ramo linear $r_t\Ahat_t$ é menor (mais negativo) e
        \textbf{mantém o gradiente ativo para desfazê-la}.
      \end{teobox}
    \end{column}
    \begin{column}{0.48\textwidth}
      \begin{alertabox}[title={Heurística vs. garantia teórica}]
        O recorte limita $r_t$ amostra a amostra, e não a KL média em esperança.\\[0.3em]
        O PPO-Clip \textbf{não garante melhoria monotônica formal}: é uma aproximação barata,
        inspirada na teoria da região de confiança.
      \end{alertabox}
    \end{column}
  \end{columns}
\end{frame}

%% ============================================================
\begin{frame}{Algoritmo: PPO com objetivo recortado}
  \framesubtitle{Arquitetura ator-crítico, como nas implementações modernas}

  \begin{algbox}{PPO-Clip}
    \small
    \textbf{Entrada:} ator $\theta_0$, crítico $w_0$, recorte $\epsilon\approx 0{,}2$, épocas $K\in[3,10]$,
    pesos $c_1\approx 0{,}5$ e $c_2\approx 0{,}01$.\\
    \textbf{para} $k = 0, 1, 2, \dots$ \textbf{faça}
    \begin{enumerate}\setlength{\itemsep}{0pt}
      \item Execute $\pol_{\theta_k}$ por $T$ passos em $N$ ambientes paralelos (lote de $NT$ transições).
      \item Calcule os alvos $\hat G_t$ e as vantagens $\Ahat_t$ (via GAE) e \textbf{congele-os}.
      \item Armazene $\log\pol_{\theta_k}(a_t\mid s_t)$ como constante: é o denominador fixo de $r_t(\theta)$.
      \item Para cada época $1,\dots,K$ e cada minilote ($\Ahat_t$ normalizada, $\|\nabla\|\le 0{,}5$),
            dê um passo de Adam na perda conjunta.
    \end{enumerate}
  \end{algbox}
  {\small\color{rlCinza} A razão é calculada em espaço log, por estabilidade numérica:}
  \[
    r_t(\theta) = \exp\big(\log\pol_\theta(a_t\mid s_t) - \log\pol_{\theta_k}(a_t\mid s_t)\big)
  \]
\end{frame}

%% ============================================================
\begin{frame}{A perda conjunta}
  \framesubtitle{Ator recortado, crítico de valor e bônus de entropia}

  \begin{teobox}{Minimizada a cada minilote}
    \[
      \mathcal{L}(\theta,w) = -L^{\mathrm{CLIP}}_{\theta_k}(\theta) + c_1\big(V_w(s_t) - \hat G_t\big)^2
        - c_2\,\mathcal{H}\big[\pol_\theta(\cdot\mid s_t)\big]
    \]
  \end{teobox}

  \begin{columns}[T]
    \begin{column}{0.32\textwidth}
      \begin{ideiabox}[title={Ator: $-L^{\mathrm{CLIP}}$}]
        Melhora a política com segurança, dentro da região recortada.
      \end{ideiabox}
    \end{column}
    \begin{column}{0.32\textwidth}
      \begin{ideiabox}[title={Crítico: $c_1(V - \hat G_t)^2$}]
        Treina o crítico de valor por regressão aos alvos congelados.
      \end{ideiabox}
    \end{column}
    \begin{column}{0.32\textwidth}
      \begin{ideiabox}[title={Exploração: $-c_2\,\mathcal{H}[\pol_\theta]$}]
        Impede o colapso prematuro para uma política determinística.
      \end{ideiabox}
    \end{column}
  \end{columns}
\end{frame}

%% ============================================================
\begin{frame}{Engenharia do PPO: três erros fatais de implementação}
  \framesubtitle{Boa parte do desempenho publicado do PPO vem destes detalhes}

  \begin{columns}[T]
    \begin{column}{0.32\textwidth}
      \begin{alertabox}[title={1 · Recalcular $\Ahat_t$ dentro do laço de $K$ épocas}]
        As vantagens medem a qualidade das ações \textbf{segundo a política que coletou os dados}
        e devem ficar congeladas durante toda a iteração $k$.
      \end{alertabox}
    \end{column}
    \begin{column}{0.32\textwidth}
      \begin{alertabox}[title={2 · Não destacar o denominador de $r_t(\theta)$}]
        $\pol_{\theta_k}(a_t\mid s_t)$ é uma constante numérica. Se o autodiff derivar também o
        denominador, $r_t(\theta)\equiv 1$ e o gradiente é \textbf{identicamente nulo}.
        Use valores armazenados ou \emph{detach}.
      \end{alertabox}
    \end{column}
    \begin{column}{0.32\textwidth}
      \begin{alertabox}[title={3 · Usar épocas $K$ em excesso}]
        O recorte por amostra não trava a distribuição inteira: com $K$ alto, $\pol_\theta$ foge
        silenciosamente de $\pol_{\theta_k}$ e o retorno degrada \textbf{sem erro explícito}.
      \end{alertabox}
    \end{column}
  \end{columns}
\end{frame}

%% ============================================================
\begin{frame}{Métricas obrigatórias no painel}
  \framesubtitle{Monitore sempre ao fim de cada iteração de treino}

  \begin{columns}[T]
    \begin{column}{0.48\textwidth}
      \begin{teobox}{Métrica 1 · KL empírica entre $\pol_\theta$ e $\pol_{\theta_k}$}
        \small Se disparar: reduza $K$, reduza a taxa de aprendizado ou aplique
        \emph{early stopping} por KL.
      \end{teobox}
    \end{column}
    \begin{column}{0.48\textwidth}
      \begin{teobox}{Métrica 2 · Fração de amostras recortadas}
        \small A fração com $|r_t - 1| > \epsilon$. Subindo, a política está mudando rápido
        demais para o lote atual.
      \end{teobox}
    \end{column}
  \end{columns}

  \begin{ideiabox}
    Ambas são baratas de calcular e denunciam o erro 3 antes que ele apareça na curva de retorno.
  \end{ideiabox}
\end{frame}

%% ============================================================
\begin{frame}{Síntese: do REINFORCE ao PPO}
  \framesubtitle{As cinco equações que conectam a aula}
  \small
  \begin{tabular}{@{}p{0.27\textwidth}l@{}}
    \toprule
    \textbf{Conceito} & \textbf{Formulação} \\
    \midrule
    1. Gradiente com linha de base &
      $\nabla_\theta J = \E\big[\sum_t\gamma^t\nabla_\theta\log\pol_\theta(a_t\mid s_t)\big(G_t-b(s_t)\big)\big],\ \
       \E\big[\nabla_\theta\log\pol_\theta\, b(s)\big]=0$ \\[0.5em]
    2. Linha de base ótima e vantagem &
      $b^*(s)\approx\Vpi(s),\quad A^{\pol} = \Qpi - \Vpi,\quad \E_{a\sim\pol}\big[A^{\pol}(s,a)\big] = 0$ \\[0.5em]
    3. Métrica local KL--Fisher &
      $\KLbar(\theta+\Delta\theta\,\|\,\theta)\approx\tfrac12\Delta\theta^\top F\,\Delta\theta,\quad
       F = \E\big[\nabla_\theta\log\pol_\theta\,\nabla_\theta\log\pol_\theta^\top\big]$ \\[0.5em]
    4. Diferença, substituto e limite &
      $J(\pol')-J(\pol) = \tfrac{1}{1-\gamma}\,\E_{d^{\pol'},\pol'}[A^{\pol}]\approx L_{\pol}(\pol'),\quad
       \text{erro}\le C\sqrt{\KLbar}$ \\[0.5em]
    5. Objetivo recortado do PPO &
      $r_t = \frac{\pol_\theta(a_t\mid s_t)}{\pol_{\theta_k}(a_t\mid s_t)},\quad
       L^{\mathrm{CLIP}} = \E\big[\sum_t\min\big(r_t\Ahat_t,\ \clipop(r_t,1-\epsilon,1+\epsilon)\,\Ahat_t\big)\big]$ \\
    \bottomrule
  \end{tabular}
\end{frame}

%% ============================================================
\begin{frame}{Exercícios fundamentais}
  \framesubtitle{Roteiro de fixação teórica e experimental}

  \begin{enumerate}\setlength{\itemsep}{0.45em}
    \item \textbf{Dependência na ação.} No bandido de 2 ações com $b(a) = r(a)$, mostre que
          $\E[\ghat] = 0 \neq \nabla_\theta J$ e aponte o passo da prova que quebra.
    \item \textbf{Invariância a deslocamento.} Com $r' = r + c$, prove que $\E[\ghat]$ não muda,
          mas $\Var[\ghat]$ cresce com $c^2$ sem linha de base.
    \item \textbf{Viés do erro de TD.} Expresse $\E[\delta_t\mid s,a] - A^{\pol}(s,a)$ em função do
          erro do crítico $\epsilon(s) = V(s;w) - \Vpi(s)$.
    \item \textbf{Substituto local.} Verifique $L_{\pol}(\pol) = 0$ e que os gradientes de $L$ e de $J$
          coincidem em $\pol'=\pol$; discuta quando $L_{\pol}(\pol') > 0$ garante $J(\pol') > J(\pol)$.
    \item \textbf{Experimento PPO.} Compare $K\in\{1,4,20\}$ épocas monitorando KL empírica,
          fração recortada e retorno médio.
  \end{enumerate}
\end{frame}

%% ============================================================
\begin{frame}{Referências essenciais}
  \begin{itemize}\setlength{\itemsep}{0.5em}
    \item \textbf{Kakade \& Langford (2002).} \emph{Approximately Optimal Approximate Reinforcement Learning.} ICML.
    \item \textbf{Schulman et al. (2015).} \emph{Trust Region Policy Optimization} (TRPO). ICML.
    \item \textbf{Mnih et al. (2016).} \emph{Asynchronous Methods for Deep Reinforcement Learning} (A3C). ICML.
    \item \textbf{Achiam, Held, Tamar \& Abbeel (2017).} \emph{Constrained Policy Optimization} (CPO). ICML.
    \item \textbf{Schulman et al. (2017).} \emph{Proximal Policy Optimization Algorithms.} arXiv:1707.06347.
  \end{itemize}

  \vspace{0.8em}
  {\small\color{rlCinza} Material adaptado de Emma Brunskill, CS234 (Stanford), e dos slides de
  Joshua Achiam sobre gradientes de política avançados.}
\end{frame}

%% ============================================================
\begin{frame}[fragile]{Anexo: o experimento de variância por estimador}
  \begin{columns}[T]
    \begin{column}{0.4\textwidth}
      \scriptsize
      \textbf{\color{rlAzul}Problema: cadeia de 8 estados}
      \begin{itemize}\setlength{\itemsep}{0pt}
        \item Estados 0 a 7, início em 0; ações: esquerda ou direita.
        \item $+1$ por passo (ruído gaussiano de desvio 2); $+40$ ao chegar ao estado 7, que encerra o episódio.
        \item Horizonte de 40 passos; $\gamma = 0{,}99$.
      \end{itemize}
      \textbf{\color{rlAzul}Agente e medição}
      \begin{itemize}\setlength{\itemsep}{0pt}
        \item Política softmax tabular, treinada com GAE ($\lambda = 0{,}95$); crítico tabular por Monte Carlo.
        \item 150 iterações de 48 episódios; os 5 pesos são avaliados no mesmo lote.
        \item Variância = traço da covariância do gradiente por episódio.
        \item Código: \texttt{aula06/sim/simulacao.py}.
      \end{itemize}
    \end{column}
    \begin{column}{0.58\textwidth}
\fontsize{6.5}{7.6}\selectfont
\begin{verbatim}
def step(s, a):              # a = 1: direita
    s2 = max(0, s + (1 if a == 1 else -1))
    if s2 == N - 1:          # chegou ao fim
        return s2, 40 + rng.normal(0, 2), True
    return s2, 1 + rng.normal(0, 2), False

for it in range(150):        # 48 episódios/iteração
    for tr in episodios:
        G   = retornos(tr, gam)
        dlt = r + gam * V[s2] - V[s]   # erro de TD
        gae = desconto(dlt, gam * lam)
        W = dict(R=G[0], G=G, GV=G - V[s],
                 GAE=gae, TD=dlt)
        # g[k] = soma_t grad log pi(a_t|s_t) * W[k][t]
        V[s] += 0.002 * (G - V[s])     # crítico MC
    var[k] = mean(norm(g[k] - mean(g[k]))**2)
    theta += 0.0005 * mean(g["GAE"])
\end{verbatim}
    \end{column}
  \end{columns}
\end{frame}

\end{document}
